\clearpage
\section*{Apartment counts for original construction and later assessments}
\textit{September 15, 2026. Agent-drafted population review requested by Jacob.}

The Madison errors expose a weakness in the earlier review: selecting an
Assessor vintage did not establish that apartment counts, floor area and
land described the same building. The 2024 source gives both 1100 and
1048 W Madison 81 apartments while preserving their individual building
and lot measurements. The 2021 counts are 9 and 24. Independent building
reports support 9 and 18, within an 81-apartment community. Both records
enter the main 500-foot sample on Burnett's side. The earlier review
accepted the newer count when permit evidence did not support the older
one, and called a positive selected count resolved. That branch did not
require evidence supporting the newer count. The active producer also
selects the latest commercial assessment without resolving these count
changes. Neither the bedroom summation nor a join created the raw 81s.

The citywide screen covers all 13,677 retained projects dated 2006--2022,
including 4,023 in the current main sample. It finds 218 projects meeting
at least one review criterion, including 60 in the main sample. These
are flags, not 218 errors. Criteria concern changes in apartment counts
on unchanged buildings and lots, historical reviews cleared without unit
evidence, development counts inconsistent with available permits,
nearby repeated counts or large lot measurements, extreme floor area per
apartment, and overlapping parcel identifiers. They were applied to the
population without using scores or estimated effects to select cases.

Eight reviewed records have supported count errors, five of them in the
main sample. Besides Madison, 1607 W Waveland is assigned eight apartments
although its original permit specifies six; an October 2023 permit
creates two more by converting its restaurant. At 922 W Washington,
the selected 149 includes 45 pre-existing lofts alongside 104 new
apartments. The developer and original permit explicitly distinguish
these components. Its floor/lot coverage and selected construction year
still require reconciliation. Conversely, 2051 N Clifton has two
apartments in the selected assessment, while the original permit,
original sale and later same-building assessment support three.

Three development assessments also understate apartment counts. At
550 N Ogden the selected source says Six, whereas the original permit
and completed condo listings report eight. At 5914 N Lincoln the
selected count is four against six in the permit and completed listings.
At 1400 W Monroe the selected count is six. Its original permit says
42, but a March 2021 revision reduces the count to 38, matching the
current developer description. The builder's completion announcement
still says 42; it should not override the explicit revision. The
recommended count is 38, with floor-area coverage and the selected 2021
year still to be reconciled. These development records prioritize a
positive Assessor apartment field over a conflicting permit count.
The preserved apartment field contains decoded words; a numeric
category-code parsing error was considered and ruled out.

A further strong conflict concerns 3408 N Milwaukee: two tied parcels
report six apartments in earlier same-building assessments and two in
later ones. The original new-building permit specifies six, but its
July 2013 issuance conflicts with the selected 2012 construction year.
This is not yet a fully resolved observation. Ten other reviewed records
have evidence supporting their current count despite a flag, including
321 W Evergreen (15) and 2109 S Halsted (18). At 3332 W Irving Park,
three original apartments should remain three even though a 2023
conversion creates a fourth. Blanket preference for either old or new
counts would therefore introduce errors.

The review records seven other unresolved measurement cases and eight
projects with large lot areas repeated on nearby records. Separate tax
parcels do not establish separate reported land coverage. In the strict
residential history comparison, 12,054 projects have a comparable card
and measurements without a count conflict, 35 have a conflict, and 778
have no comparable history. Passing these screens does not independently
validate every building, year or area. The remaining flagged population
has not received complete individual adjudication.

No production corrections or estimates were changed. The next data work
must distinguish the original construction from later conversions,
resolve the consequential count and area conflicts, and apply adopted
corrections once in the ordinary cleaning path. The magnitude of the
resulting regression change is unknown because no re-estimation was run.

\textit{Reproduction.} The audit is
\path{tasks/audits/construction_measurement_review/}. Its two scripts
produce the population screen, exact commercial source rows, full-date
permit candidates and HTML evidence report; \path{review_notes.csv}
contains 34 documented findings or unresolved questions. Source links
and permit identifiers appear there and in the README. The audit uses
preserved production outputs following commit \texttt{32c094b8} on
\texttt{score\_robustness}; Make held those outputs fixed to avoid a
production rebuild during review. The saved analysis snapshot matches
the current analysis file byte for byte, and all its observations match
the citywide ledger on units, floor area, land and year. These are
checks of the reviewed inputs, not a fresh production replication.

\paragraph{Coverage assessment before more manual review.}
Jacob asked whether 218 represents the universe of potential problems.
It does not. Of the flagged records, 126 have a source disagreement or an
earlier count review cleared without supporting evidence; 92 trigger
only the nearby-repeat or unusual-ratio checks. Sixty-four already have
a recorded decision and 146 have an earlier commercial count review;
these groups overlap. Their evidence should be reused for the particular
variable it addresses. None of the eight confirmed count-error projects
has a case-specific entry in the current correction table. These errors
point to gaps in ordinary source selection and its automated review,
rather than evidence that all previous manual adjudications were wrong.

The current screens leave 771 unflagged residential projects without a
comparable historical assessment and 113 unflagged commercial projects
without a positive older count. They also leave 62 commercial projects
unflagged despite changes in floor or lot measurements, because area
changes alone are not a comprehensive criterion. Fifty development
records lack an available positive permit count. These groups overlap;
the counts describe missing comparison coverage, not new errors. Stable
but incorrect source values can pass, and the retained-project screen
cannot establish that construction was not incorrectly omitted.

Before further individual adjudication, the checks should cover all
required measurements and relevant source branches, reuse the existing
decisions, and be assessed against an independently selected sample of
unflagged records. Catching the eight known errors does not validate the
screen's miss rate, because some checks were added after seeing those
cases. No such independent validation has yet been completed. The
coverage assessment changes no production values, estimates or flags.
